\chapter{Paths and nodes}
\label{ch:paths}

Sooner or later an imported drawing is almost right, or a shape you need does not
exist in any library. That is when you work below the level of ``rectangle and
circle'' and edit the actual geometry. This chapter covers the drawing tools,
the node editor, and the operations that reshape paths.

\section{What an element actually is}

Elements come in a few kinds, and the kind decides which tools apply:

\begin{center}\small
\begin{tabular}{@{}L{3.0cm}L{11.2cm}@{}}
\toprule
\textbf{Kind} & \textbf{Behaviour} \\
\midrule
Rectangle, ellipse, circle, line, polyline & \emph{Parametric} shapes: they
remember their definition, so width, height, radius and corner behaviour stay
editable. \\
Path & Free geometry: a sequence of points, lines and curves. Imported SVG
drawings are almost always paths. \\
Text & Either bitmap text (a rendered glyph block) or vector text from a font
file (Chapter~\ref{ch:text}). \\
Image & A bitmap (Chapter~\ref{ch:images}). \\
Group, layer, file & Containers. They have no geometry of their own. \\
\bottomrule
\end{tabular}
\end{center}

\noindent \menu{Convert to path} turns a parametric shape or a text element into a
path, after which it can be edited node by node but no longer by its original
parameters. Two siblings of that command exist for undoing the transformation in
spirit: \btn{Reify user changes} commits the changes you made by dragging (so
that later parameter edits no longer fight with them), and
\btn{Reset user changes} throws those changes away and restores the parametric
shape.

\begin{tipbox}[Convert late]
Keep shapes parametric while you are still deciding sizes --- a rectangle you
can still set to 60\unit{mm} in the \panel{Position} pane is worth more than a
path you have to measure. Convert when you need to move individual corners, not
before.
\end{tipbox}

\section{Drawing}

The drawing tools live in the \panel{Tools} icon bar and each has simple click
semantics:

\begin{center}\small
\begin{tabular}{@{}L{2.6cm}L{11.6cm}@{}}
\toprule
\textbf{Tool} & \textbf{How it works} \\
\midrule
\btn{Line} & Click the start, click the end. \\
\btn{Rectangle}, \btn{Circle}, \btn{Ellipse} & Click one corner or centre, click
the opposite. Set exact sizes afterwards in the \panel{Position} pane. \\
\btn{Polygon}, \btn{Polyline} & Left click adds a point and a segment;
double-click completes the shape; right-click cancels.
\emph{Polygon} offers sub-modes: \emph{Freehand}, \emph{Regular},
\emph{Star 1}, \emph{Star 2} and \emph{Grow}. \\
\btn{Vector} & The bezier tool: click for a corner point, click and hold to pull
out a curve handle, double-click to complete, right-click to end an open path.
This is the tool to reach for when you are tracing something by hand. \\
\btn{Draw} & Freehand drawing: the path follows your pointer. \\
\btn{Point} & Places a point element --- the input for a \mkseries{Dots}
operation (drilling, perforation). \\
\btn{Text} & Places a text element. \\
\btn{Job Start} & Places a job-starting position on the scene
(\emph{placement}). \\
\bottomrule
\end{tabular}
\end{center}

\noindent Two ribbon modes are worth knowing because they are not shapes at all:
\btn{Measure} draws dimension lines (the overlays visible in
Figure~\ref{fig:loaded}) without adding elements to the job, and the
\emph{nudge} ribbon mode lets you push a drawn line around as if it were physical
--- used for organic, hand-drawn-looking paths.

\begin{notebox}[Drawing is not design]
Everything in this chapter works, and people do build real parts with it. If
your artwork is more than a dozen curves, draw it in a vector editor and import
it. Chapter~\ref{ch:files} covers the round trip, including the trick of
right-clicking a file node and choosing \btn{Reload} after editing the source.
\end{notebox}

\section{Editing nodes}

Select a path and activate \btn{Node Edit} in the ribbon. The tool shows the
path's nodes --- the points where segments meet and where curves are controlled
--- and lets you work on them directly:

\begin{itemize}
  \item \textbf{Drag a node} to move it. Snapping applies, so a node can be
        dropped precisely onto a grid intersection or another element's point.
  \item \textbf{\btn{Node Move}} is the simpler sibling tool, intended purely for
        dragging nodes to new positions.
  \item \textbf{\btn{Point Move}} does the same job for point elements and
        placements.
  \item \textbf{\btn{Parametric Edit}} appears for parametric shapes: instead of
        editing raw nodes you edit the shape's underlying parameters, so a star
        stays a star while you change its point count or inner radius.
  \item \textbf{\btn{Tab Edit}} edits the \emph{tabs} of an element --- the short
        uncut segments that hold a part in place during cutting, so it does not
        shift or fall through. Tabs are what make a cut sheet survive the last
        few seconds of the job.
\end{itemize}

\noindent Node editing is the one place where the geometry is genuinely fragile:
a path that looks closed may have a 0.02\unit{mm} gap, and a laser will happily
leave a part hanging by that gap. If a shape should be a closed loop, check it
afterwards by cutting it on scrap.

\section{Reshaping paths}

These commands live on the element context menu, or in the console:

\begin{center}\small
\begin{tabular}{@{}L{4.9cm}L{9.3cm}@{}}
\toprule
\textbf{Command} & \textbf{What it does} \\
\midrule
\btn{Break Subpaths} & Splits one path element into its separate subpaths, so
each can be moved, assigned or deleted on its own. The cure for a drawing where
``one object'' is really forty. \\
\btn{Merge elements} & Joins selected paths into a single element. \\
\btn{Group elements} / \btn{Ungroup elements} & Groups are organisational only:
they do not alter geometry and do not merge paths. \\
\btn{Simplify group} & Reduces nested group structure without changing geometry. \\
\btn{Remove transparent objects} & Deletes elements that would burn nothing (no
stroke and no fill). Useful after a heavy import. \\
\btn{Outline element(s)} & Creates a new path offset by a given distance --- for
example \cmd{1mm}. Not the same as kerf compensation: this makes a new element
rather than adjusting the operation. \\
\menu{Offset shapes} & \btn{\dots to outside} and \btn{\dots to inside} by a
chosen distance (1--5\unit{mm}), used for lids, masks and press fits. \\
\cmd{stitch [tolerance]} & Connects path ends that are within the tolerance of
each other, reducing the number of separate burns and the travel between them.
Console only; the cut planner can also stitch automatically
(Chapter~\ref{ch:order}). \\
\cmd{clipper} & Boolean combination of selected elements: the first argument
is \cmd{union}, \cmd{difference}, \cmd{intersection} or \cmd{xor}. Console only --- there is no menu entry for it --- with
options for fill type, interpolation and whether the originals are kept. \\
\cmd{offset <distance>}, \cmd{offset2}, \cmd{offsetsimple} & The console
equivalents of the offset menus, with options for join type, separate output and
interpolation. \\
\cmd{pocket <offset> <repeats>} & Pocketing: repeated inward offsets that hollow
an area out, one ring per repeat. \\
\bottomrule
\end{tabular}
\end{center}

\section{Generating shapes you cannot draw}

Three generators produce geometry that would be tedious to draw by hand.

\subsection*{Living hinges}

The \btn{Hinge} window generates a bendable hinge pattern --- the classic
laser-cut lattice that lets plywood flex. You give it the \emph{X} and \emph{Y}
coordinate of the hinge area, its \emph{Width} and \emph{Height}, and choose a
\emph{Pattern} and \emph{Algorithm} from the combos, with a \emph{Rotation} and
the \emph{Cell-Width} / \emph{Cell-Height} ratios that control the lattice.
\btn{Preview Pattern} shows the result, \btn{Preview Shape} shows the boundary it
will occupy, \btn{Apply to scene} puts it into the drawing, and
\btn{Default Values} resets the numbers. Living hinges are cut, not engraved, so
assign the result to a \mkseries{Cut} operation and expect a slow job: this is a
lot of path.

\subsection*{Tabs and keyholes}

\btn{Tab Edit} places the small uncut bridges described above. The keyhole tool
(right-click an element, \btn{Add a keyhole}) does the opposite job: it combines
two elements so that one is engraved through the other as a mask --- used for
nested artwork and for engraving a picture into a shape.

\subsection*{Parametric generators (console)}

A family of console commands builds mathematical shapes directly:

\begin{lstlisting}[style=konsole]
pgrid 10mm 10mm 4 3          # a grid of copies: cell size, columns, rows
shape 5 0 0 20mm             # star/polygon: corners, x, y, radius
growingshape 2 2 5 6 10mm    # a growing spiral-like shape
cycloid 2 2 10mm 4mm 20      # cycloid curve
fractal_tree 2 2 5 6         # a fractal tree
\end{lstlisting}

\noindent Each takes sizes and counts; \cmd{shape} and \cmd{growingshape} accept
options such as \cmd{--angle}, \cmd{--ratio} and \cmd{--density}. They are not in
any menu, which is deliberate: they are the escape hatch for the shapes a GUI
cannot anticipate. Once created, the result is a normal element that the
\btn{Parametric Edit} tool can continue to adjust.

\begin{tryit}{Rescue a messy import}
Take an SVG that has too many objects and make it cuttable.
\begin{enumerate}[label=\arabic*.,leftmargin=1.4em]
  \item Import the file and expand the tree. Note how many elements arrive and
        whether they are grouped.
  \item Use \btn{Measure} to check one known dimension against the real artwork.
        If it is wrong, fix the size in the \panel{Position} pane before anything
        else.
  \item Find a path that should be several objects, select it, and use
        \btn{Break Subpaths}. Move the pieces apart to see what they are.
  \item Use \btn{Remove transparent objects} to shed anything invisible.
  \item Add a \mkseries{Cut} operation for the outlines and a \mkseries{Engrave}
        operation for the interior detail, assigning by colour or by selection.
  \item Use \cmd{stitch 0.1mm} in the console and compare the plan's travel
        distance before and after in the \panel{Simulation} window.
  \item Save as a new project file. Your rescue is now reusable artwork.
\end{enumerate}
\end{tryit}
